\documentclass[../../../include/open-logic-section]{subfiles}

\begin{document}
\olfileid{sth}{ord-arithmetic}{intro}

\olsection{ভূমিকা}

\olref[ordinals][]{chap}-এ আমরা ক্রমসংখ্যার একটি তত্ত্ব গড়েছি।
\olref[spine][]{chap}-এ দেখেছি, ক্রমসংখ্যাগুলিকে এমন একটি
মেরুদণ্ড ভাবা যায়, যার চারপাশে স্তরক্রমের অবশিষ্ট অংশ গড়ে ওঠে।
কিন্তু ক্রমসংখ্যার ভূমিকা শুধু এটুকুই নয়।
ক্রমসংখ্যার পাটিগণিতও করতে হয়।

এর ইঙ্গিত আমরা আগেই দিয়েছি।
\olref[ordinals][idea]{sec}-এ $\omega$,
$\omega+1$ এবং $\omega+\omega$-র কথা বলার সময়
আমাদের আলোচনা ছিল অনানুষ্ঠানিক; এবার তা যথাযথভাবে
ব্যাখ্যা করার পালা। তবে উল্লেখ করা ভালো, এই অধ্যায়ে
দর্শনের আলোচনা বেশি নেই। এখানে মূলত কারিগরি বিকাশ
থাকবে, আর সেই সঙ্গে একটি কিছুটা আকর্ষণীয় পর্যবেক্ষণ:
একই কাজ দুটি ভিন্ন উপায়ে করা যায়।

\end{document}
